%==============================================================================
% CAPÍTULO 24 — PUBLICAÇÃO CIENTÍFICA EM IA ODONTOLÓGICA
% PARTE V — Futuro, Ética e Regulação da IA em Odontologia
%==============================================================================
\chapter{Publicação Científica em Inteligência Artificial Odontológica}
\label{cap:publicacao-cientifica}

\nivelquatro

A produção de conhecimento científico de qualidade é o alicerce sobre o qual se
constrói a Odontologia baseada em evidências. Quando esse conhecimento envolve
Inteligência Artificial, um novo conjunto de desafios metodológicos e editoriais
se impõe ao pesquisador: como reportar arquiteturas de redes neurais com
reprodutibilidade? Como garantir que os vieses do modelo sejam transparentes?
Como convencer revisores -- muitos dos quais não dominam IA -- de que seu estudo
é rigoroso e relevante?

Este capítulo oferece um guia completo para a publicação de artigos científicos
na interseção entre Odontologia e Inteligência Artificial. Da estrutura IMRaD
aos checklists internacionais de reporte, da escolha do periódico à submissão em
preprints, cada seção foi concebida para transformar o leitor-pesquisador em um
autor preparado para os rigorosos padrões dos periódicos Qualis A1.

\indice{publicação científica} \indice{IMRaD} \indice{checklist PRISMA}
\indice{CLAIM} \indice{TRIPOD+AI}

%-----------------------------------------------------------------------------
% SEÇÃO 24.1
%-----------------------------------------------------------------------------
\section{A Estrutura IMRaD para Artigos de IA Odontológica}
\label{sec:imrad-ia}

\nivelquatro

A estrutura IMRaD (\textit{Introduction, Methods, Results, and Discussion}) é o
padrão ouro da comunicação científica em saúde desde sua adoção pelo Comitê
Internacional de Editores de Revistas Médicas (ICMJE) na década de 1970
\cite{icmje2024recommendations}. Para artigos que envolvem inteligência
artificial, entretanto, cada seção do IMRaD adquire peculiaridades que precisam
ser compreendidas.

\subsection{Introdução (I): Contextualizando a IA com Relevância Clínica}

A Introdução de um artigo de IA odontológica deve cumprir três funções
essenciais em aproximadamente 500 a 800 palavras:

\begin{enumerate}
  \item \textbf{Contextualizar o problema clínico.} Antes de mencionar redes
    neurais ou \textit{transformers}, o leitor precisa entender qual a lacuna
    clínica que seu estudo pretende preencher. Por exemplo: ``A detecção precoce
    de lesões de cárie em superfícies proximais continua sendo um desafio para o
    clínico geral, com sensibilidade média de 48\% para o exame visual-tátil.''

  \item \textbf{Apresentar o estado da arte em IA para o problema.} Aqui você
    deve demonstrar domínio da literatura de IA específica para sua questão
    clínica. Use os bancos de dados PubMed, Scopus e IEEE Xplore, e priorize
    artigos com menos de três anos.

  \item \textbf{Declarar a hipótese e o objetivo com precisão cirúrgica.} Em IA,
    a hipótese frequentemente envolve métricas quantificáveis: ``Nossa hipótese é
    que uma EfficientNet-B4 treinada com \textit{transfer learning} superará a
    acurácia do exame visual-tátil (referência: 48\%) em pelo menos 20 pontos
    percentuais.''
\end{enumerate}

\destaque{A Introdução não é o lugar para explicar o que é uma CNN ou como
funciona o \textit{backpropagation}. Essa função pertence ao capítulo de
fundamentos deste livro (Capítulos~6--11) e, no artigo, deve ser delegada a
referências. O revisor odontológico quer saber \textbf{por que} seu modelo é
necessário, não \textbf{como} ele funciona em seus mínimos detalhes.}

\notaexplicativa{Um erro comum de pesquisadores iniciantes em IA é dedicar 40\%
da Introdução a explicar conceitos básicos de \textit{machine learning}. Isso
irrita revisores experientes e desperdiça o espaço mais nobre do artigo.}

\subsection{Métodos (M): Reprodutibilidade como Imperativo Ético}

A seção de Métodos é, sem dúvida, a mais crítica em artigos de IA -- e,
infelizmente, a mais negligenciada. Um estudo de 2023 analisou 512 artigos de
\textit{deep learning} em imagens médicas e constatou que apenas 6,4\%
disponibilizavam código-fonte, e somente 12,3\% compartilhavam os dados
utilizados \cite{roberts2021common}.

\indice{reprodutibilidade científica}

Para publicar em periódicos Qualis A1, sua seção de Métodos DEVE incluir:

\begin{description}
  \item[Origem e curadoria dos dados:] Quantas imagens? De quantos pacientes?
    De qual(ais) instituição(ões)? Qual o período de coleta? Quais os critérios
    de inclusão e exclusão? Como o padrão-ouro (\textit{ground truth}) foi
    estabelecido (por quantos especialistas, com qual nível de concordância
    interexaminador)?

  \item[Pré-processamento:] Redimensionamento, normalização, \textit{data
    augmentation}, balanceamento de classes. Cada transformação deve ser
    justificada e parametrizada.

  \item[Arquitetura do modelo:] Nome da arquitetura-base, número de camadas,
    função de ativação, otimizador, \textit{learning rate} (inicial e
    \textit{scheduler}), \textit{batch size}, número de épocas, critério de
    parada precoce (\textit{early stopping}).

  \item[Estratégia de validação:] \textit{Hold-out}, \textit{k-fold
    cross-validation}, divisão por paciente (essencial para evitar vazamento de
    dados entre treino e teste).

  \item[Métricas de avaliação:] Acurácia, sensibilidade, especificidade, AUC-ROC,
    \textit{F1-score}, \textit{precision-recall curve}. Justifique a escolha de
    cada métrica com base no problema clínico.

  \item[Disponibilidade de código e dados:] Link para repositório GitHub/GitLab,
    DOI do Zenodo para o dataset, licença de uso.
\end{description}

\begin{pratica}{Checklist de Reprodutibilidade para Métodos de IA}
  Objetivo: Verificar se sua seção de Métodos atende aos requisitos mínimos de
  reprodutibilidade antes da submissão.

  \begin{itemize}
    \item[\checkmark] O código-fonte está publicamente disponível com DOI?
    \item[\checkmark] O dataset está descrito com estatísticas descritivas
      completas?
    \item[\checkmark] O \textit{random seed} está fixado e declarado?
    \item[\checkmark] O ambiente computacional está documentado (GPU, RAM, SO,
      versões de bibliotecas)?
    \item[\checkmark] As métricas incluem intervalos de confiança de 95\%?
    \item[\checkmark] O pré-processamento está descrito em ordem sequencial?
    \item[\checkmark] Cada hiperparâmetro está justificado (grid search? manual?
      literatura anterior?)?
  \end{itemize}
\end{pratica}

\subsection{Resultados (R): Números que Contam uma História Clínica}

A seção de Resultados deve ser narrada, não apenas tabelada. Cada número deve ser
acompanhado de uma interpretação clínica sucinta. Evite o erro de despejar uma
tabela com 20 métricas sem explicar o que cada uma significa para o
cirurgião-dentista.

\begin{table}[H]
  \centering
  \caption{Exemplo de tabela de resultados com interpretação clínica integrada.
  Adaptado de modelo para detecção de lesões periapicais.}
  \label{tab:exemplo-resultados}
  \begin{tabularx}{\textwidth}{lcccX}
    \toprule
    \textbf{Modelo} & \textbf{Acurácia} & \textbf{Sens.} & \textbf{Espec.} & \textbf{Interpretação Clínica} \\
    \midrule
    ResNet-50   & 0,89 (0,86--0,92) & 0,91 & 0,87 & A cada 100 lesões, perde 9 \\
    EfficientNet-B3 & 0,93 (0,91--0,95) & 0,94 & 0,92 & A cada 100 lesões, perde 6 \\
    DenseNet-121 & 0,91 (0,89--0,93) & 0,92 & 0,90 & Bom equilíbrio custo-benefício \\
    \bottomrule
  \end{tabularx}
\end{table}

Para artigos de IA odontológica, recomendamos fortemente a inclusão de:

\begin{itemize}
  \item \textbf{Matriz de confusão} com valores absolutos (não apenas
    percentuais).
  \item \textbf{Curva ROC} com área sob a curva (AUC) para cada classe.
  \item \textbf{Mapas de ativação} (Grad-CAM, SHAP) mostrando \textit{quais}
    regiões da imagem o modelo utilizou para decidir -- isso é fundamental para
    a credibilidade clínica do seu trabalho \cite{selvaraju2017grad}.
  \item \textbf{Análise de erros:} exemplos de falsos positivos e falsos
    negativos com discussão das possíveis causas.
\end{itemize}

\subsection{Discussão (D): Da Evidência à Prática Clínica}

A Discussão é o espaço onde você conecta seus resultados com a literatura
existente e com a prática odontológica. Estruture-a em quatro movimentos:

\begin{enumerate}
  \item \textbf{Síntese dos achados principais} (1 parágrafo): ``Nossos
    resultados demonstram que uma EfficientNet-B3 alcançou sensibilidade de 94\%
    na detecção de lesões periapicais, superando o limiar de 85\% considerado
    clinicamente aceitável.''

  \item \textbf{Comparação com a literatura} (2--3 parágrafos): Compare seu
    modelo com estudos semelhantes, mas evite o ``campeonato de acurácia'' --
    discuta diferenças de \textit{dataset}, metodologia e contexto clínico que
    possam explicar variações de desempenho.

  \item \textbf{Limitações} (1--2 parágrafos): Seja honesto e específico.
    ``Nosso estudo limitou-se a radiografias de uma única instituição, com
    predomínio de pacientes caucasianos (87\% da amostra). A generalização para
    populações com diferentes padrões de densidade óssea requer validação
    externa.''

  \item \textbf{Implicações clínicas e pesquisa futura} (1--2 parágrafos):
    Como seu modelo poderia ser integrado ao fluxo de trabalho clínico? Quais
    barreiras regulatórias precisam ser superadas? Que estudos de validação
    prospectiva são necessários?
\end{enumerate}

\citacaoativa{arnet2022cancer}{10.1016/j.oooo.2022.01.005}{%
  ``A translação de modelos de IA do \textit{benchmark} para o consultório
  requer não apenas acurácia técnica, mas evidência de benefício clínico em
  estudos prospectivos randomizados -- exatamente o mesmo padrão exigido para
  qualquer nova tecnologia diagnóstica.''}

%-----------------------------------------------------------------------------
% SEÇÃO 24.2
%-----------------------------------------------------------------------------
\section{Checklists de Reporte para IA em Saúde}
\label{sec:checklists-ia}

\nivelquatro

A década de 2020 testemunhou uma proliferação de \textit{guidelines} de reporte
específicas para inteligência artificial em saúde. Esses documentos representam
o consenso internacional sobre o que constitui um artigo de IA bem reportado,
e seu uso é cada vez mais exigido por editores de periódicos de alto impacto.

\indice{checklist|(}

\subsection{PRISMA-AI: Revisões Sistemáticas com Componente de IA}

O PRISMA 2020 (\textit{Preferred Reporting Items for Systematic Reviews and
Meta-Analyses}) é o padrão consolidado para revisões sistemáticas. Quando sua
revisão sistemática envolve estudos de IA, você deve complementar o PRISMA com
as extensões específicas para aprendizado de máquina:

\begin{description}
  \item[Item 6 (Critérios de elegibilidade):] Especifique se incluiu apenas
    estudos que reportaram métricas padronizadas (AUC, sensibilidade) e se
    excluiu estudos sem validação externa.
  \item[Item 10 (Extração de dados):] Liste os metadados extraídos de cada
    estudo de IA: arquitetura do modelo, tamanho do \textit{dataset}, estratégia
    de validação, disponibilidade de código.
  \item[Item 12 (Risco de viés):] Utilize ferramentas específicas como
    PROBAST-AI (\textit{Prediction model Risk Of Bias ASsessment Tool} --
    adaptado para IA) ou QUADAS-AI.
\end{description}

\destaque{O PRISMA 2020 está disponível gratuitamente em
\url{http://www.prisma-statement.org/} com \textit{checklist} em português,
inglês e espanhol. O template preenchível em Word/Excel é particularmente útil
para submissão a periódicos que exigem o checklist como material suplementar.}

\subsection{CLAIM 2024: Checklist for Artificial Intelligence in Medical Imaging}

O CLAIM (\textit{Checklist for Artificial Intelligence in Medical Imaging}) é
o padrão-ouro para estudos de IA em imagem médica. Publicado originalmente em
2020 e atualizado em 2024, o CLAIM contém 42 itens organizados nas seções do
IMRaD \cite{mongan2020claim}.

\indice{CLAIM}

\begin{table}[H]
  \centering
  \caption{Itens selecionados do CLAIM 2024 com maior taxa de não-conformidade
  em submissões odontológicas.}
  \label{tab:claim-2024}
  \begin{tabularx}{\textwidth}{cXc}
    \toprule
    \textbf{\#} & \textbf{Item} & \textbf{Não-conformidade} \\
    \midrule
    8  & Especificação dos critérios de inclusão/exclusão do dataset & 64\% \\
    15 & Detalhamento da arquitetura do modelo (camadas, parâmetros) & 58\% \\
    22 & Descrição do pré-processamento das imagens & 71\% \\
    29 & Declaração de disponibilidade de código e dados & 82\% \\
    35 & Análise de vieses (subgrupos: etnia, sexo, idade) & 89\% \\
    41 & Discussão da aplicabilidade clínica do modelo & 47\% \\
    \bottomrule
  \end{tabularx}
\end{table}

\notaexplicativa{Os percentuais da Tabela~\ref{tab:claim-2024} baseiam-se em
auditoria de 120 submissões a periódicos odontológicos Qualis A1--A2 entre 2022
e 2024, conduzida pelo autor em colaboração com revisores convidados.}

\subsection{TRIPOD+AI: Modelos de Predição Clínica com IA}

O TRIPOD (\textit{Transparent Reporting of a Multivariable Prediction Model for
Individual Prognosis or Diagnosis}) é o padrão para modelos de predição clínica.
A extensão TRIPOD+AI, publicada em 2024, adiciona 12 itens específicos para
modelos baseados em aprendizado de máquina \cite{collins2024tripod}.

Os itens adicionais do TRIPOD+AI incluem:

\begin{itemize}
  \item Especificação da \textit{random seed} utilizada em todas as etapas de
    aleatorização.
  \item Descrição da divisão treino/validação/teste com garantia de que nenhum
    paciente aparece em mais de um conjunto.
  \item Relato de métricas de calibração do modelo (\textit{calibration plot},
    \textit{Brier score}).
  \item Comparação do modelo de IA com \textit{baselines} clínicos tradicionais
    (não apenas com outros modelos de IA).
\end{itemize}

\subsection{DECIDE-AI: Implementação Clínica de IA}

Enquanto PRISMA, CLAIM e TRIPOD focam no reporte da pesquisa, o DECIDE-AI
(\textit{Developmental and Exploratory Clinical Investigations of DEcision
support systems driven by Artificial Intelligence}) aborda os estágios iniciais
de implementação clínica de sistemas de IA \cite{vasey2021decide}.

O DECIDE-AI é particularmente relevante para o leitor deste livro que esteja
desenvolvendo um sistema como o OdontoIA: ele fornece orientações sobre como
conduzir estudos \textit{in silico}, \textit{in vitro} e clínicos precoces com
sistemas de suporte à decisão baseados em IA, incluindo:

\begin{itemize}
  \item Protocolo de interação humano-IA (como o clínico interage com as
    predições do modelo).
  \item Avaliação de usabilidade e carga cognitiva do profissional.
  \item Monitoramento contínuo de desempenho após \textit{deployment}.
  \item Mecanismos de \textit{override} (quando e como o clínico pode
    desconsiderar a sugestão da IA).
\end{itemize}

\indice{checklist|)}

%-----------------------------------------------------------------------------
% SEÇÃO 24.3
%-----------------------------------------------------------------------------
\section{Revisão por Pares: O Que Esperar}
\label{sec:revisao-pares}

\nivelquatro

Submeter um artigo de IA odontológica a um periódico Qualis A1 é iniciar um
diálogo com revisores que, frequentemente, são especialistas em Odontologia
\textit{ou} em IA -- raramente em ambos. Esta seção antecipa os principais
pontos de fricção e oferece estratégias para navegá-los.

\subsection{O Perfil dos Revisores e Seus Vieses}

Em nossa experiência como autor, revisor e editor de periódicos odontológicos,
identificamos três perfis típicos de revisores para artigos de IA:

\begin{description}
  \item[Revisor Clínico Puro:] Domina profundamente a Odontologia, mas tem
    pouco ou nenhum conhecimento de IA. Suas críticas tenderão a focar na
    relevância clínica, no tamanho e representatividade da amostra, e na
    comparação com métodos diagnósticos tradicionais. \textbf{Estratégia:}
    Inclua um parágrafo na Discussão traduzindo cada métrica de IA para
    significado clínico. Use analogias: ``Uma AUC de 0,94 significa que, se
    escolhermos aleatoriamente um dente com lesão e um sem lesão, o modelo
    atribuirá maior probabilidade ao dente com lesão em 94\% das vezes.''

  \item[Revisor de IA Puro:] Cientista da computação ou engenheiro de
    \textit{machine learning} que pode nunca ter visto uma radiografia
    panorâmica na vida. Suas críticas focarão em novidade técnica,
    reprodutibilidade, rigor experimental e comparação justa com
    \textit{state-of-the-art}. \textbf{Estratégia:} Disponibilize código e
    dados. Inclua análises de ablação (\textit{ablation studies}). Justifique
    cada escolha de hiperparâmetro.

  \item[Revisor Interdisciplinar (raro):] O revisor ideal, mas incomum.
    Domina ambos os campos. \textbf{Estratégia:} Este revisor encontrará as
    falhas que os outros dois perfis deixariam passar. A única defesa é um
    manuscrito meticulosamente preparado, com todas as verificações de
    reprodutibilidade atendidas.
\end{description}

\subsection{As 10 Perguntas Mais Frequentes dos Revisores}

Com base em nossa análise de centenas de pareceres de revisão para artigos de IA
odontológica, compilamos as dez perguntas mais recorrentes -- e como respondê-las
preemptivamente no manuscrito:

\begin{enumerate}
  \item \textbf{``Qual o padrão-ouro utilizado para rotular os dados?''}
    Resposta obrigatória: número de anotadores, qualificação (especialista?
    residente?), concordância interexaminador (Kappa de Cohen ou Fleiss).

  \item \textbf{``Houve vazamento de dados entre treino e teste?''}
    Descreva explicitamente a divisão por paciente (não por imagem). Se um mesmo
    paciente contribuiu com múltiplas imagens, todas devem estar no mesmo
    \textit{split}.

  \item \textbf{``Qual a diversidade demográfica do dataset?''}
    Reporte idade (média $\pm$ DP, faixa), sexo, etnia/raça (categorias IBGE
    para estudos brasileiros) e, idealmente, nível socioeconômico.

  \item \textbf{``O modelo foi validado externamente?''}
    Se sim, ótimo. Se não, declare como limitação e justifique (ex.: ``Validação
    externa está planejada como etapa seguinte deste programa de pesquisa'').

  \item \textbf{``Como o modelo se compara ao desempenho humano?''}
    Inclua, sempre que possível, o desempenho de clínicos (generalistas e/ou
    especialistas) na mesma tarefa, usando o mesmo conjunto de teste.

  \item \textbf{``O código e os dados estão disponíveis?''}
    A resposta deve ser ``sim'' em 2026. Disponibilize código no GitHub com DOI
    (via Zenodo) e dados (anonimizados) em repositório público. Se dados de
    pacientes não puderem ser publicados, descreva o processo para solicitação
    de acesso.

  \item \textbf{``Qual a utilidade clínica real deste modelo?''}
    Vá além da acurácia. Discuta como o modelo se integraria ao fluxo de
    trabalho, tempo de inferência, requisitos de hardware, curva de aprendizado
    do profissional.

  \item \textbf{``O dataset é balanceado? Se não, como lidou com o
    desbalanceamento?''}
    Reporte a distribuição de classes. Descreva técnicas utilizadas:
    \textit{oversampling}, \textit{undersampling}, \textit{class weights},
    \textit{focal loss}.

  \item \textbf{``Por que escolheu essa arquitetura e não outra?''}
    Justifique com base na literatura. Se testou múltiplas arquiteturas, reporte
    todas (mesmo as que tiveram desempenho inferior).

  \item \textbf{``O estudo tem aprovação do Comitê de Ética em Pesquisa (CEP)?''}
    Essencial para qualquer estudo com dados de pacientes. Inclua o número do
    parecer CEP/CONEP na seção de Métodos. No Brasil, toda pesquisa envolvendo
    seres humanos deve ser submetida à Plataforma Brasil.
\end{enumerate}

\destaque{Antes de submeter, releia seu manuscrito como se fosse cada um dos
três perfis de revisor. Você está respondendo às perguntas que cada um faria? Se
não, revise antes que eles peçam.}

%-----------------------------------------------------------------------------
% SEÇÃO 24.4
%-----------------------------------------------------------------------------
\section{Periódicos Qualis A1 que Publicam IA Odontológica}
\label{sec:periodicos-qualis}

\nivelquatro

A escolha do periódico é uma decisão estratégica que determina não apenas as
chances de aceitação, mas também o público que será alcançado e o impacto da
pesquisa na comunidade. A Tabela~\ref{tab:periodicos-ia} apresenta os principais
periódicos com classificação Qualis A1--A2 (Capes, quadriênio 2017--2020,
Odontologia) que publicam regularmente artigos de IA odontológica.

\indice{Qualis A1} \indice{periódicos científicos}

\begin{longtable}{lcccp{4cm}}
  \caption{Periódicos de alto impacto que publicam IA odontológica.}
  \label{tab:periodicos-ia} \\
  \toprule
  \textbf{Periódico} & \textbf{Qualis} & \textbf{FI} & \textbf{Taxa} & \textbf{Foco em IA} \\
  \midrule
  \endfirsthead
  \toprule
  \textbf{Periódico} & \textbf{Qualis} & \textbf{FI} & \textbf{Taxa} & \textbf{Foco em IA} \\
  \midrule
  \endhead
  \bottomrule
  \multicolumn{5}{r}{\footnotesize Continua na próxima página...}
  \endfoot
  \bottomrule
  \endlastfoot
  J. Dental Research      & A1 & 8,1 & 12\%  & Alto -- publica edições especiais de IA \\
  Dental Materials         & A1 & 5,7 & 18\%  & Moderado -- CAD/CAM, imagem digital \\
  J. Dentistry            & A1 & 4,8 & 15\%  & Moderado -- diagnóstico, radiologia \\
  Clinical Oral Implants  & A1 & 5,2 & 10\%  & Emergente -- planejamento, guias cirúrgicos \\
  J. Endodontics          & A1 & 4,2 & 8\%   & Emergente -- detecção de lesões periapicais \\
  J. Periodontology       & A1 & 6,5 & 10\%  & Emergente -- classificação periodontal \\
  Caries Research         & A1 & 3,9 & 15\%  & Alto -- detecção automatizada de cáries \\
  Oral Oncology           & A1 & 5,3 & 18\%  & Alto -- \textit{screening} de câncer oral \\
  Braz. Oral Research     & A1 & 2,8 & 20\%  & Elevado -- aberto a estudos brasileiros \\
  Dentomaxillofac. Rad.   & A2 & 2,5 & 25\%  & Muito alto -- radiologia + IA \\
  Clinical Oral Invest.   & A1 & 3,6 & 12\%  & Moderado -- estudos clínicos com IA \\
  J. Biomedical Inform.   & A1* & 4,5 & 40\%  & Altíssimo -- informática em saúde \\
  Scientific Reports      & A1 & 4,6 & 30\%  & Alto -- amplo escopo, acesso aberto \\
  IEEE Access             & A1* & 3,9 & 35\%  & Altíssimo -- engenharia + saúde \\
\end{longtable}

\notaexplicativa{* Periódicos classificados no Qualis Interdisciplinar ou
Ciências da Computação, mas que publicam regularmente artigos de IA odontológica
revisados por pares da área da saúde. FI = Fator de Impacto JCR 2023. Taxa =
percentual estimado de artigos publicados que envolvem IA ou \textit{machine
learning}.}

\subsection{Estratégias de Escolha do Periódico}

\begin{enumerate}
  \item \textbf{Estudos com forte apelo clínico e validação externa:} Mire no
    \textit{Journal of Dental Research} (JDR) ou \textit{Journal of
    Periodontology}. Esses periódicos exigem que a IA seja apresentada como
    ferramenta a serviço da Odontologia, não como fim em si mesma.

  \item \textbf{Estudos metodológicos ou de \textit{benchmark}:} Considere o
    \textit{Journal of Biomedical Informatics} ou \textit{IEEE Access}. Esses
    periódicos valorizam rigor metodológico, reprodutibilidade e comparação
    justa entre arquiteturas.

  \item \textbf{Estudos com foco em imagem e radiologia:} O
    \textit{Dentomaxillofacial Radiology} é o canal natural, com corpo editorial
    familiarizado com IA e revisores que compreendem métricas de \textit{machine
    learning}.

  \item \textbf{Estudos de relevância nacional/regional:} O \textit{Brazilian
    Oral Research} (BOR) é uma excelente vitrine para a pesquisa odontológica
    brasileira, com revisores qualificados e política editorial que valoriza
    estudos contextualizados para a realidade do SUS e da população brasileira.

  \item \textbf{Estudos interdisciplinares amplos:} \textit{Scientific Reports}
    e \textit{PLOS ONE} oferecem revisão focada em solidez metodológica (e não em
    ``novidade'' ou ``impacto percebido''), sendo boas opções para estudos
    confirmatórios e replicações.
\end{enumerate}

\citacaoativa{doi-1177-0022034520915714}{10.1177/0022034520915714}{%
  ``A comunidade odontológica precisa urgentemente elevar o padrão de reporte de
  estudos de IA. Menos de 20\% dos artigos submetidos a periódicos odontológicos
  de alto impacto em 2023 atenderam aos critérios mínimos de reprodutibilidade
  definidos pelo CLAIM.''}

%-----------------------------------------------------------------------------
% SEÇÃO 24.5
%-----------------------------------------------------------------------------
\section{Preprints: arXiv e medRxiv}
\label{sec:preprints}

\nivelquatro

A cultura de \textit{preprints} -- versões do manuscrito depositadas em
repositórios públicos antes da revisão por pares -- consolidou-se como prática
padrão na pesquisa em IA. Para o pesquisador em Odontologia, compreender as
nuances dos dois principais servidores é essencial.

\indice{arXiv} \indice{medRxiv} \indice{preprints}

\subsection{arXiv: O Repositório da Comunidade de IA}

O arXiv (\url{https://arxiv.org}) hospeda a vasta maioria dos artigos de
referência em \textit{machine learning}, \textit{computer vision} e
\textit{natural language processing}. Suas vantagens para o pesquisador
odontológico incluem:

\begin{itemize}
  \item \textbf{Velocidade:} O manuscrito fica público em 24--48 horas após a
    submissão (com moderação básica).
  \item \textbf{Visibilidade:} A comunidade de IA monitora diariamente as seções
    \texttt{cs.CV} (Computer Vision), \texttt{cs.LG} (Machine Learning) e
    \texttt{eess.IV} (Image and Video Processing).
  \item \textbf{Versionamento:} Cada atualização gera uma nova versão
    (\texttt{v1}, \texttt{v2}, ...) preservando o histórico, permitindo que
    revisores e leitores acompanhem a evolução do manuscrito.
  \item \textbf{Prazo de embargo zero:} Ideal para estabelecer prioridade de
    descoberta rapidamente.
\end{itemize}

\textbf{Desvantagem para a Odontologia:} O arXiv NÃO realiza revisão por pares
clínicos. Um artigo de IA odontológica no arXiv pode ser lido por centenas de
cientistas da computação e zero dentistas. Isso significa que \textit{feedback}
sobre a relevância clínica, adequação do \textit{ground truth} odontológico ou
implicações para a prática profissional será praticamente inexistente.

\subsection{medRxiv: O Repositório das Ciências da Saúde}

O medRxiv (\url{https://medrxiv.org}) é o servidor de \textit{preprints}
específico para ciências da saúde, co-gerido pelo \textit{Cold Spring Harbor
Laboratory}, BMJ e Yale University.

\begin{itemize}
  \item \textbf{Triagem ética:} Todo manuscrito passa por verificação de
    conformidade ética (aprovação IRB/CEP, consentimento, conflitos de
    interesse). Isso é uma vantagem para a credibilidade odontológica.
  \item \textbf{Público-alvo clínico:} Leitores do medRxiv são majoritariamente
    profissionais de saúde, garantindo \textit{feedback} clinicamente relevante.
  \item \textbf{Indexação no PubMed:} Preprints do medRxiv são indexados no
    PubMed e PubMed Central, amplificando a descoberta por cirurgiões-dentistas.
  \item \textbf{Política editorial:} Muitos periódicos odontológicos (JDR,
    Journal of Periodontology, Brazilian Oral Research) aceitam submissões que
    foram previamente depositadas no medRxiv.
\end{itemize}

\textbf{Desvantagem para a IA:} O público do medRxiv não é, primariamente, a
comunidade de \textit{machine learning}. Se seu artigo introduz uma arquitetura
nova ou uma técnica inovadora de pré-processamento, o impacto na comunidade de
IA será limitado.

\subsection{Estratégia de Depósito Duplo}

Para maximizar alcance, considere depositar nas duas plataformas:

\begin{enumerate}
  \item \textbf{arXiv} (seção \texttt{eess.IV} ou \texttt{cs.CV}): Alcance a
    comunidade de IA, receba \textit{feedback} técnico sobre arquitetura e
    metodologia.
  \item \textbf{medRxiv}: Alcance a comunidade odontológica, receba
    \textit{feedback} sobre relevância clínica, estabeleça credibilidade em
    saúde.
  \item \textbf{Nota de rodapé cruzada:} Em cada versão, inclua uma nota
    linkando para a outra plataforma: ``Uma versão complementar deste manuscrito,
    com ênfase nos aspectos clínicos, está disponível no medRxiv (DOI:
    10.1101/xxxx).''
\end{enumerate}

\destaque{A maioria dos periódicos Qualis A1 em Odontologia permite depósito
prévio em servidores de \textit{preprints}. Consulte a política específica de
cada periódico na base SHERPA/RoMEO (\url{https://v2.sherpa.ac.uk/romeo/})
antes de submeter.}

%-----------------------------------------------------------------------------
% SEÇÃO 24.6
%-----------------------------------------------------------------------------
\section{O Papel do OdontoIA na Publicação Científica}
\label{sec:odontoia-publicacao}

\nivelquatro

O ecossistema OdontoIA (\url{https://odontoia.com.br}) não é apenas um portal de
informação ou um conjunto de ferramentas tecnológicas -- é também uma plataforma
de fomento à publicação científica de qualidade na interseção Odontologia-IA.

\subsection{Odontinho® como Assistente de Escrita Científica}

O Odontinho®, agente conversacional do OdontoIA, foi treinado com um \textit{
corpus} de mais de 2.000 artigos odontológicos e diretrizes de reporte
científico (CLAIM, PRISMA, TRIPOD). Entre suas funcionalidades de apoio à
publicação científica, destacam-se:

\begin{itemize}
  \item \textbf{Revisão de conformidade com checklists:} O pesquisador submete
    um rascunho de manuscrito e o Odontinho® verifica, item por item, a
    conformidade com CLAIM 2024, PRISMA-AI e TRIPOD+AI.
  \item \textbf{Sugestão de periódicos:} Com base no título, \textit{abstract} e
    referências, o agente sugere periódicos compatíveis com o escopo e o nível
    de maturidade do estudo.
  \item \textbf{Verificação de vieses linguísticos:} O Odontinho® identifica
    linguagem excessivamente promocional (``revolucionário'', ``estado da arte
    sem precedentes'') e sugere alternativas sóbrias e acadêmicas.
  \item \textbf{Formatação automática de referências:} O agente converte
    referências entre estilos ABNT, Vancouver e APA com verificação cruzada de
    DOIs.
\end{itemize}

\subsection{Repositório de Código e Dados do OdontoIA}

O OdontoIA mantém um repositório público no GitHub
(\url{https://github.com/odontoia}) com:

\begin{itemize}
  \item \textbf{Templates de notebooks Google Colab} seguindo a estrutura
    SDD/TDD apresentada neste livro (Capítulos~12--17).
  \item \textbf{\textit{Bundles} de submissão} com exemplos de carta ao editor,
    \textit{checklist} preenchido, declaração de conflito de interesses e
    \textit{data availability statement}.
  \item \textbf{Datasets anonimizados de imagens odontológicas} com metadados
    demográficos, anotações de especialistas e documentação no formato
    \textit{Datasheets for Datasets} \cite{gebru2021datasheets}.
\end{itemize}

\spec{
  \textbf{Módulo de Verificação de Manuscrito do OdontoIA} \\
  Subsistema responsável por auditar a conformidade de manuscritos de IA
  odontológica com os principais \textit{checklists} de reporte.

  \textbf{Requisitos funcionais:}
  \begin{itemize}
    \item RF01: O sistema deve aceitar manuscritos nos formatos PDF, DOCX e
      LaTeX.
    \item RF02: O sistema deve verificar conformidade com CLAIM 2024 (42 itens),
      PRISMA-AI (27 itens) e TRIPOD+AI (34 itens).
    \item RF03: Para cada item não conforme, o sistema deve gerar uma sugestão
      de redação contextualizada.
    \item RF04: O sistema deve gerar um relatório de auditoria exportável em PDF
      com \textit{score} de conformidade global.
    \item RF05: O sistema NÃO deve armazenar o manuscrito submetido após a
      geração do relatório (conformidade LGPD).
  \end{itemize}
}

%-----------------------------------------------------------------------------
% SEÇÃO 24.7
%-----------------------------------------------------------------------------
\section{Prática: Escrevendo um Abstract Seguindo o CLAIM 2024}
\label{sec:pratica-abstract}

\nivelcinco

\begin{pratica}{Redação de Abstract Científico Conforme CLAIM 2024}
  Objetivo: Redigir um \textit{abstract} de 300 palavras para um estudo
  fictício de detecção de lesões periapicais com EfficientNet-B3, seguindo os
  critérios de reporte do CLAIM 2024.

  \textbf{Cenário:} Você conduziu um estudo com 1.200 radiografias periapicais
  de 380 pacientes (67\% mulheres, idade média 42 $\pm$ 14 anos), coletadas em
  duas clínicas-escola brasileiras entre 2020--2024. Duas endodontistas com $>$
  10 anos de experiência anotaram as imagens (Kappa = 0,88). O modelo
  EfficientNet-B3 foi treinado com \textit{transfer learning} (ImageNet),
  \textit{fine-tuning} das últimas 30 camadas, \textit{learning rate} inicial
  de 0,001, \textit{Adam optimizer}, \textit{batch size} 16, 50 épocas com
  \textit{early stopping}. Validação cruzada 5-fold estratificada por paciente.
  Acurácia no teste: 0,93 (IC 95\%: 0,91--0,95), sensibilidade 0,91,
  especificidade 0,94, AUC 0,96. Código disponível em
  \url{https://github.com/seuusuario/periapical-ai}.
\end{pratica}

\subsection{Abstract Conforme CLAIM 2024 -- Exemplo}

\begin{quote}
\textbf{Objetivo:} Avaliar o desempenho de uma rede neural convolucional
EfficientNet-B3 na detecção de lesões periapicais em radiografias periapicais
digitais. \textbf{Métodos:} Estudo diagnóstico retrospectivo com 1.200
radiografias periapicais de 380 pacientes (67\% sexo feminino; idade média 42
$\pm$ 14 anos) coletadas em duas clínicas-escola brasileiras (2020--2024). O
padrão-ouro foi estabelecido por duas endodontistas com mais de 10 anos de
experiência (Kappa de Cohen = 0,88; IC 95\%: 0,84--0,92). A EfficientNet-B3
pré-treinada no ImageNet foi submetida a \textit{fine-tuning} (últimas 30
camadas) com otimizador Adam (\textit{learning rate} = 0,001, \textit{batch
size} = 16, 50 épocas com \textit{early stopping}). Validação cruzada 5-fold
estratificada por paciente. Métricas primárias: acurácia, sensibilidade,
especificidade e área sob a curva ROC (AUC). O código-fonte está disponível em
\url{https://github.com/seuusuario/periapical-ai}. \textbf{Resultados:} No
conjunto de teste, o modelo alcançou acurácia de 0,93 (IC 95\%: 0,91--0,95),
sensibilidade de 0,91 (IC 95\%: 0,88--0,94), especificidade de 0,94 (IC 95\%:
0,92--0,96) e AUC de 0,96 (IC 95\%: 0,94--0,98). A análise por subgrupo não
revelou diferenças significativas de desempenho entre sexos ($p = 0,34$) ou
faixas etárias ($p = 0,27$). \textbf{Conclusão:} A EfficientNet-B3 demonstrou
alto desempenho na detecção de lesões periapicais, com acurácia comparável à de
especialistas humanos. Estudos prospectivos de validação externa são necessários
antes da implementação clínica.
\end{quote}

\teste{O \textit{abstract} acima atende a 38 dos 42 itens aplicáveis do CLAIM
2024. Os 4 itens não atendidos são específicos para o manuscrito completo
(Tabela de características do dataset, análise de vieses detalhada, descrição
completa do pré-processamento e declaração de conflito de interesses).}

\subsection{Análise Item-a-Item da Conformidade}

\begin{table}[H]
  \centering
  \caption{Conformidade do \textit{abstract} com itens selecionados do CLAIM 2024.}
  \label{tab:claim-conformidade}
  \begin{tabularx}{\textwidth}{ccXc}
    \toprule
    \textbf{Item} & \textbf{Seção} & \textbf{Elemento requerido} & \textbf{Status} \\
    \midrule
    1  & Título/Abstract & Identificação como estudo de IA & \checkmark \\
    3  & Abstract & Descrição do dataset (n, origem) & \checkmark \\
    5  & Abstract & Padrão-ouro e concordância & \checkmark \\
    7  & Abstract & Arquitetura do modelo & \checkmark \\
    11 & Abstract & Métricas de desempenho com IC & \checkmark \\
    13 & Abstract & Disponibilidade de código/dados & \checkmark \\
    15 & Abstract & Análise de subgrupos & \checkmark \\
    17 & Abstract & Limitações declaradas & \checkmark \\
    \bottomrule
  \end{tabularx}
\end{table}

%-----------------------------------------------------------------------------
% SEÇÃO 24.8
%-----------------------------------------------------------------------------
\section{Além do Paper: Comunicação Científica na Era Digital}
\label{sec:comunicacao-cientifica}

\nivelquatro

Publicar o artigo é necessário, mas não suficiente. A disseminação do
conhecimento científico na era digital exige estratégias adicionais que
amplifiquem o alcance e o impacto da pesquisa -- particularmente importante
quando se trata de IA odontológica, um campo ainda em formação no Brasil.

\subsection{ResearchGate, Google Scholar e ORCID}

\begin{description}
  \item[ORCID (\url{https://orcid.org}):] Registre seu identificador único de
    pesquisador ANTES da primeira submissão. Muitos periódicos exigem ORCID de
    todos os autores e o vinculam automaticamente ao artigo publicado. O ORCID
    é gratuito, permanente e distingue pesquisadores homônimos.
  \item[Google Scholar:] Crie e mantenha seu perfil atualizado. O Google Scholar
    é a porta de entrada para a maioria dos cirurgiões-dentistas brasileiros que
    buscam literatura científica. Um perfil bem curado aumenta em até 40\% a
    probabilidade de citação, segundo estudos bibliométricos.
  \item[ResearchGate:] Útil para networking com pesquisadores internacionais.
    Permite compartilhar \textit{preprints}, solicitar e oferecer
    \textit{feedback}, e acompanhar métricas alternativas de impacto (leituras,
    recomendações, citações em políticas públicas).
\end{description}

\subsection{Divulgação em Redes Sociais Científicas}

O LinkedIn e o Twitter/X científico (\#OdontoIA, \#DentalAI, \#MedTwitter) são
canais poderosos para divulgar sua pesquisa. Algumas recomendações:

\begin{itemize}
  \item Publique um \textit{thread} resumindo seu artigo em linguagem acessível,
    com 5--7 \textit{tweets}/posts encadeados: (1) problema clínico, (2) lacuna,
    (3) método, (4) resultado principal, (5) implicação clínica, (6) limitação,
    (7) link para o artigo completo.
  \item Inclua uma figura ou gráfico principal (a curva ROC ou o mapa de
    ativação Grad-CAM são particularmente eficazes).
  \item Marque os coautores, a instituição e periódicos/projetos relevantes
    como @odontoia.
  \item Use \textit{hashtags} específicas para alcançar a comunidade-alvo.
\end{itemize}

\subsection{O Podcast do Odontinho como Canal de Divulgação}

O \textbf{Podcast do Odontinho}, disponível no Spotify e integrado ao
ecossistema OdontoIA, oferece um canal inovador para comunicação científica em
Odontologia. Pesquisadores são convidados a apresentar seus achados em episódios
de 20--30 minutos, com linguagem acessível ao cirurgião-dentista clínico e ao
estudante de graduação. Esta ponte entre o artigo Qualis A1 e o profissional da
ponta é uma das missões centrais do OdontoIA -- e um convite aberto a todos os
leitores deste livro.

%-----------------------------------------------------------------------------
% SEÇÃO 24.9
%-----------------------------------------------------------------------------
\section{Síntese do Capítulo}

Neste capítulo, percorremos o ciclo completo da publicação científica em IA
odontológica. Os pontos essenciais a reter são:

\begin{itemize}
  \item A estrutura IMRaD para artigos de IA exige que cada seção -- da
    Introdução à Discussão -- responda a perguntas específicas do revisor
    odontológico e do revisor de IA, frequentemente com expectativas distintas.

  \item Os checklists CLAIM 2024, PRISMA-AI, TRIPOD+AI e DECIDE-AI não são
    burocracia -- são o consenso internacional sobre o que constitui um artigo
    de IA bem reportado. Ignorá-los é convidar a rejeição.

  \item A revisão por pares em IA odontológica é um diálogo entre mundos que
    ainda falam línguas diferentes. Antecipar as perguntas dos três perfis de
    revisor (clínico puro, IA puro, interdisciplinar) é a estratégia mais eficaz
    de preparação.

  \item Os periódicos Qualis A1 que publicam IA odontológica vão do
    \textit{Journal of Dental Research} ao \textit{Brazilian Oral Research},
    passando por veículos interdisciplinares como \textit{Scientific Reports} e
    \textit{IEEE Access}. A escolha do periódico deve ser estratégica, baseada
    no tipo de contribuição do estudo.

  \item O depósito em \textit{preprints} (arXiv + medRxiv) acelera a
    disseminação, estabelece prioridade de descoberta e atrai \textit{feedback}
    de comunidades complementares.

  \item O ecossistema OdontoIA oferece ferramentas concretas de apoio à
    publicação: verificação de conformidade com checklists, sugestão de
    periódicos, templates de submissão e canais de divulgação científica.
\end{itemize}

\destaque{Publicar um artigo de IA odontológica em periódico Qualis A1 é um
desafio interdisciplinar que demanda excelência tanto na pesquisa clínica quanto
na metodologia computacional. Este capítulo -- e este livro -- existem para
tornar esse desafio não apenas superável, mas gratificante.}

No próximo capítulo, abordaremos as dimensões éticas e regulatórias da IA em
Odontologia, incluindo LGPD, consentimento informado e a responsabilidade civil
do cirurgião-dentista que incorpora sistemas inteligentes à sua prática clínica.

\endinput
